% ============================================================
% DS MOD08 -- SUJET B : loi binomiale, loi de Poisson, loi normale
% BTS CPRP 1re année -- durée 1 h 30
% Pôle Formation UIMM - CVDL - S. JAUBERT
% Corrigé et barème : MOD08_DS_Binomiale_Poisson_Normale_CPRP_Sujet_B_Corrige_v1_0 (sujet jumeau : Sujet_A)
% ============================================================
\documentclass[a4paper,11pt]{article}
\usepackage[utf8]{inputenc}
\usepackage[T1]{fontenc}
% babel-french absent sur certains postes : chargement conditionnel
\IfFileExists{french.ldf}{\usepackage[french]{babel}}{}
\usepackage{helvet}
\usepackage{courier}
\renewcommand{\familydefault}{\sfdefault}
\usepackage[top=3cm, bottom=2.5cm, left=2cm, right=2cm, headheight=2cm]{geometry}
\usepackage{amsmath,amssymb,amsfonts}
\usepackage{graphicx}
\usepackage[dvipsnames,table]{xcolor}
\usepackage{tikz}
\usepackage[most]{tcolorbox}
\usepackage{fancyhdr}
\usepackage{enumitem}
\usepackage{array,needspace}
\usepackage{titlesec}
\usepackage{hyperref}

\definecolor{uimmblue}{HTML}{005C85}
\definecolor{uimmred}{HTML}{D31326}
\definecolor{uimmlightblue}{HTML}{E8F1F8}
\definecolor{uimmgray}{HTML}{333333}
\definecolor{uimmlightgray}{HTML}{F5F5F5}
\hypersetup{colorlinks=true, linkcolor=uimmblue, urlcolor=uimmblue}

\titleformat{\section}{\Large\bfseries\color{uimmblue}}{}{0pt}{}[\vspace{2pt}{\color{uimmblue}\titlerule[1.2pt]}]
\setlist[enumerate,1]{label={\color{uimmblue}\textbf{\arabic*.}}, leftmargin=*}
\setlist[enumerate,2]{label=\alph*), leftmargin=*}
\setlist[itemize,1]{label={\color{uimmblue}\textbullet}}

\newtcolorbox{exo}[1]{colback=white, colframe=uimmblue, fonttitle=\bfseries\color{white},
  coltitle=white, title={#1}, boxrule=1pt, arc=2pt,
  left=8pt, right=8pt, top=5pt, bottom=6pt, before skip=10pt, after skip=6pt}
\newtcolorbox{consigne}{colback=uimmlightgray, colframe=uimmgray, boxrule=0.8pt, arc=3pt,
  left=8pt, right=8pt, top=6pt, bottom=6pt}

\newcommand{\pts}[1]{\hfill\mbox{\color{uimmred}\small\textbf{(#1)}}}
% Zone de réponse lignée : #1 = nombre de lignes
\newcommand{\zone}[1]{\par\smallskip\noindent\begin{tikzpicture}
  \foreach \i in {1,...,#1}{\draw[uimmgray!30] (0,-0.85*\i) -- (\linewidth-1pt,-0.85*\i);}
  \end{tikzpicture}\par\medskip}
\newcommand{\EnTete}[2]{%
\pagestyle{fancy}\fancyhf{}\renewcommand{\headrulewidth}{0pt}
\fancyhead[L]{\raisebox{-0.5\height}{\includegraphics[height=1.5cm]{logo_uimm_placeholder.jpg}}}
\fancyhead[C]{\begin{minipage}{7cm}\centering
  {\small\color{uimmblue}\textbf{Pôle Formation UIMM - CVDL}}\\
  {\footnotesize\color{uimmgray} BTS CPRP 1\textsuperscript{re} année\\ Mathématiques}\end{minipage}}
\fancyhead[R]{{\small\color{uimmred}\textbf{MOD08 -- #1}}}
\fancyfoot[C]{{\color{uimmblue}\rule{\textwidth}{1pt}}\\[2pt]
  {\footnotesize\color{uimmgray} Pôle Formation UIMM -- Centre-Val de Loire \hfill S. JAUBERT \hfill #2 -- \thepage}}}
\EnTete{DS Sujet B}{Sujet B -- DS 1 h 30}
\newtcolorbox{rappel}{colback=uimmlightblue, colframe=uimmblue, fonttitle=\bfseries\color{white},
  coltitle=white, title={Rappels de cours}, boxrule=1pt, arc=2pt,
  left=8pt, right=8pt, top=5pt, bottom=6pt, before skip=8pt, after skip=8pt}

\begin{document}

\begin{tcolorbox}[colback=uimmred, colframe=uimmred, sharp corners, boxrule=0pt,
  left=14pt, right=14pt, top=10pt, bottom=10pt]
{\color{white}\LARGE\bfseries Devoir surveillé -- MOD08 \hfill Sujet B}\\[4pt]
{\color{white}\large Loi binomiale, loi de Poisson, loi normale \hfill Durée~: 1 h 30}
\end{tcolorbox}

\vspace{2pt}
\noindent{\small\begin{tabular}{@{}>{\bfseries\color{uimmblue}}l@{\hspace{8pt}}p{13.8cm}@{}}
Nom, prénom~: & \dotfill \hspace{1cm}{\bfseries\color{uimmblue}Note~:}\ \dots\dots\ / 20 \\
\end{tabular}}

\begin{consigne}
\textbf{Consignes.} Calculatrice autorisée en mode examen, aucun document. Répondre sur le sujet, dans les zones prévues.
\begin{itemize}[leftmargin=12pt, itemsep=1pt, topsep=2pt]
  \item Toute loi utilisée doit être \textbf{nommée avec ses paramètres}.
  \item Pour chaque probabilité demandée~: écrire l'évènement à l'aide de la variable aléatoire, \textbf{puis} le résultat. Un résultat seul ne rapporte pas tous les points.
  \item Probabilités arrondies à $10^{-4}$, cotes arrondies au micron ($10^{-3}$\,mm). Le barème est donné à titre indicatif.
\end{itemize}
\end{consigne}

\begin{rappel}
\textbf{Loi binomiale $\mathcal{B}(n\,;p)$.} On répète $n$ fois, de façon identique et indépendante, une épreuve à deux issues (succès de probabilité $p$, échec de probabilité $1-p$). La variable $X$ qui compte le nombre de succès suit la loi $\mathcal{B}(n\,;p)$. Elle prend les valeurs $0, 1, \dots, n$.
\[
P(X = k) = \binom{n}{k}\,p^{k}\,(1-p)^{n-k},\quad \binom{n}{k} = \frac{n!}{k!\,(n-k)!},\quad E(X) = np,\quad \sigma(X) = \sqrt{np(1-p)}.
\]
\textbf{Loi de Poisson $\mathcal{P}(\lambda)$.} Elle modélise un nombre d'évènements rares (pannes, défauts, arrivées) sur une durée ou une zone donnée. Elle prend les valeurs $0, 1, 2, \dots$
\[
P(X = k) = e^{-\lambda}\,\frac{\lambda^{k}}{k!},\qquad E(X) = \lambda,\qquad \sigma(X) = \sqrt{\lambda}.
\]
\textbf{Loi normale $\mathcal{N}(\mu\,;\sigma)$.} Le second paramètre est l'\textbf{écart-type} $\sigma$. La variable $Z = \dfrac{X-\mu}{\sigma}$ suit la loi $\mathcal{N}(0\,;1)$.\\
Règle des 68 -- 95 -- 99,7~: $P(\mu-\sigma\leq X\leq\mu+\sigma)\approx 0{,}68$, \ $P(\mu-2\sigma\leq X\leq\mu+2\sigma)\approx 0{,}95$, \ $P(\mu-3\sigma\leq X\leq\mu+3\sigma)\approx 0{,}997$.

\smallskip
\textbf{Calculatrice.} Lois discrètes~: «~Fdp~» (ou «~pd~») donne $P(X = k)$, «~FRép~» (ou «~cd~») donne $P(X \leq k)$.\\
{\small TI~: \texttt{binomFdp(n,p,k)}, \texttt{binomFRép(n,p,k)}, \texttt{poissonFdp($\lambda$,k)}, \texttt{poissonFRép($\lambda$,k)}. \quad Casio~: \texttt{Bpd}, \texttt{Bcd}, \texttt{Ppd}, \texttt{Pcd}.}\\
Loi normale~: $P(a\leq X\leq b)$ et valeur $x$ telle que $P(X\leq x) = p$.\\
{\small TI~: \texttt{normalFRép(a,b,$\mu$,$\sigma$)}, \texttt{FracNormale(p,$\mu$,$\sigma$)}. \quad Casio~: \texttt{Ncd}, \texttt{InvN}.}
\end{rappel}

% ------------------------------------------------------------
\newpage
\section{Exercice 1 -- Tournage de bagues \hfill {\normalsize\color{uimmred}7 points}}
\begin{exo}{Contexte}
Un tour à commande numérique usine des bagues en acier. Le suivi qualité montre que $5\,\%$ des bagues présentent un état de surface \textbf{non conforme} (rugosité $Ra$ hors spécification au rugosimètre). Le contrôleur prélève $25$ bagues dans la production. La production est assez importante pour assimiler ce prélèvement à un tirage avec remise.\\
On note $X$ le nombre de bagues \textbf{non conformes} dans le prélèvement.
\end{exo}
\begin{enumerate}
  \needspace{6.5cm}
  \item Justifier que $X$ suit une loi binomiale. Préciser ses paramètres. \pts{1,5 pt}
  \zone{5}
  \needspace{5.6cm}
  \item Calculer l'espérance $E(X)$ et l'écart-type $\sigma(X)$. Interpréter $E(X)$ dans le contexte. \pts{1,5 pt}
  \zone{4}
  \needspace{5.6cm}
  \item Calculer la probabilité que le prélèvement contienne exactement $2$ bagues non conformes. Détailler le calcul à l'aide de la formule, \textbf{sans} utiliser les fonctions de loi de la calculatrice. \pts{1 pt}
  \zone{4}
  \needspace{4cm}
  \item Pour chacun des évènements suivants, traduire l'évènement à l'aide de $X$, puis calculer sa probabilité.
  \begin{enumerate}
    \needspace{4cm}
    \item Le prélèvement contient au moins une bague non conforme. \pts{1 pt}
    \zone{3}
    \needspace{4.8cm}
    \item Au moins $22$ bagues du prélèvement sont conformes. \pts{1 pt}
    \zone{3}
    \needspace{4.8cm}
    \item Le prélèvement contient plus de $2$ bagues non conformes. \pts{1 pt}
    \zone{3}
  \end{enumerate}
\end{enumerate}

% ------------------------------------------------------------
\newpage
\section{Exercice 2 -- Casses de fil en électroérosion \hfill {\normalsize\color{uimmred}6 points}}
\begin{exo}{Contexte}
Un atelier d'outillage découpe des matrices sur une machine d'\textbf{électroérosion à fil}. Le fil se casse parfois en cours d'usinage (encrassement, mauvais arrosage, variation d'épaisseur). Ces casses surviennent de façon aléatoire et indépendante. Le suivi de production montre qu'il se produit \textbf{en moyenne $2{,}8$ casses de fil par semaine}. On admet que le nombre de casses de fil au cours d'une semaine suit une loi de Poisson.
\end{exo}
\begin{enumerate}
  \needspace{4.8cm}
  \item Définir une variable aléatoire adaptée à la situation. Donner sa loi et justifier la valeur de son paramètre. \pts{1 pt}
  \zone{3}
  \needspace{3.2cm}
  \item Calculer l'écart-type de cette variable aléatoire. \pts{0,5 pt}
  \zone{2}
  \needspace{5.6cm}
  \item Calculer la probabilité qu'il ne se produise aucune casse dans la semaine, puis la probabilité qu'il se produise exactement $3$ casses. Détailler les calculs à l'aide de la formule. \pts{1 pt}
  \zone{4}
  \needspace{4.8cm}
  \item Calculer la probabilité qu'il se produise au moins $5$ casses dans la semaine. \pts{1 pt}
  \zone{3}
  \needspace{7.4cm}
  \item Chaque casse demande l'intervention du régleur. Le chef d'atelier veut réserver dans le planning du régleur le temps nécessaire pour traiter au plus $k$ casses par semaine. Déterminer la plus petite valeur de $k$ telle que la probabilité que le nombre de casses de la semaine ne dépasse pas $k$ soit au moins égale à $0{,}90$. Justifier par deux probabilités. \pts{1,5 pt}
  \zone{6}
  \needspace{5.6cm}
  \item On admet que le nombre de casses sur une période de \textbf{trois semaines} suit une loi de Poisson dont le paramètre est le nombre moyen de casses sur trois semaines. Calculer la probabilité qu'il se produise au plus $5$ casses sur trois semaines. \pts{1 pt}
  \zone{4}
\end{enumerate}

% ------------------------------------------------------------
\newpage
\section{Exercice 3 -- Tournage d'axes de transmission \hfill {\normalsize\color{uimmred}7 points}}
\begin{exo}{Contexte}
Un tour à commande numérique usine des axes de transmission. La cote fonctionnelle est le diamètre de portée $\varnothing\,40$\,g7~: un axe est \textbf{conforme} si son diamètre est compris entre $39{,}966$\,mm et $39{,}991$\,mm.\\
Un axe trop gros (diamètre supérieur à $39{,}991$\,mm) peut être \textbf{repris}. Un axe trop petit (diamètre inférieur à $39{,}966$\,mm) est \textbf{mis au rebut}.\\
On note $D$ le diamètre, en mm, d'un axe prélevé au hasard dans la production. On admet que $D$ suit la loi normale $\mathcal{N}(39{,}976\,;0{,}005)$.
\end{exo}
\begin{enumerate}
  \needspace{4.8cm}
  \item Calculer la probabilité qu'un axe soit conforme. \pts{1 pt}
  \zone{3}
  \needspace{6.5cm}
  \item Calculer la proportion d'axes mis au rebut, puis la proportion d'axes à reprendre. \pts{1,5 pt}
  \zone{5}
  \needspace{5.6cm}
  \item À l'aide de la règle des 68 -- 95 -- 99,7, donner l'intervalle qui contient environ $99{,}7\,\%$ des diamètres. Comparer cet intervalle à la tolérance et indiquer quel côté de la tolérance pose problème. \pts{1 pt}
  \zone{4}
  \needspace{6.5cm}
  \item Le contrôleur fixe une \textbf{cote d'alerte} $a$~: $97\,\%$ des axes ont un diamètre supérieur à $a$. Traduire cette condition par une probabilité, puis calculer $a$. \pts{1,5 pt}
  \zone{5}
  \needspace{4cm}
  \item Pour réduire le rebut, le régleur décale la moyenne $\mu$ du diamètre, sans modifier l'écart-type $0{,}005$\,mm. Il vise une proportion d'axes mis au rebut d'environ $0{,}1\,\%$.
  \begin{enumerate}
    \needspace{6.5cm}
    \item Déterminer la valeur de $\mu$ à régler, arrondie au micron. \pts{1,25 pt}
    \zone{5}
    \needspace{5.6cm}
    \item Le régleur règle $\mu = 39{,}981$\,mm. Calculer la nouvelle proportion d'axes à reprendre. Le réglage est-il intéressant~? Justifier en une phrase. \pts{0,75 pt}
    \zone{4}
  \end{enumerate}
\end{enumerate}

\end{document}
